\section{Componentes biconexas por vértices}
\whenuse{Separar un grafo no dirigido en bloques que no se desconectan al borrar
un solo vértice; construir el block-cut tree o localizar articulaciones.}

Una componente biconexa por vértices es un bloque maximal que no puede separarse
eliminando uno de sus vértices. Los puntos de articulación pueden pertenecer a
varios bloques; una arista puente forma por sí sola un bloque de dos vértices.

Por ejemplo:

\begin{verbatim}
0---1
 \ /
  2---3
      / \
     4---5
\end{verbatim}

produce los bloques
$\{0,1,2\}$, $\{2,3\}$ y $\{3,4,5\}$. Los vértices $2$ y $3$ son
articulaciones porque aparecen en más de un bloque.

Durante el DFS se apilan las aristas recorridas. Si para un hijo \texttt{to} se
cumple
\[
    \texttt{low[to] >= discovery[vertex]},
\]
su subárbol no puede alcanzar un antecesor de \texttt{vertex} sin pasar por él;
las aristas de la pila hasta \texttt{vertex--to} forman un bloque completo.

\complexity{Tiempo $O(n+m)$; memoria $O(n+m)$.}

\begin{lstlisting}[style=cpp]
struct VertexBiconnectedComponents {
    struct StackedEdge { int from, to, id; };

    int timer = 0;
    vector<int> discovery, low;
    vector<StackedEdge> edge_stack;
    vector<vector<int>> components;
    const vector<vector<pair<int,int>>> &graph; // (to, edge_id)

    explicit VertexBiconnectedComponents(
            const vector<vector<pair<int,int>>> &input_graph)
        : discovery(input_graph.size(), -1),
          low(input_graph.size()), graph(input_graph) {
        for (int vertex = 0; vertex < (int)graph.size(); ++vertex) {
            if (discovery[vertex] != -1) continue;
            if (graph[vertex].empty()) {
                discovery[vertex] = low[vertex] = timer++;
                components.push_back({vertex});
            } else {
                dfs(vertex, -1);
            }
        }
    }

    void extract_component(int stop_edge) {
        vector<int> vertices;
        while (true) {
            auto edge = edge_stack.back();
            edge_stack.pop_back();
            vertices.push_back(edge.from);
            vertices.push_back(edge.to);
            if (edge.id == stop_edge) break;
        }
        sort(vertices.begin(), vertices.end());
        vertices.erase(unique(vertices.begin(), vertices.end()),
                       vertices.end());
        components.push_back(move(vertices));
    }

    void dfs(int vertex, int parent_edge) {
        discovery[vertex] = low[vertex] = timer++;

        for (auto [to, edge_id] : graph[vertex]) {
            if (edge_id == parent_edge) continue;

            if (discovery[to] == -1) {
                edge_stack.push_back({vertex, to, edge_id});
                dfs(to, edge_id);
                low[vertex] = min(low[vertex], low[to]);

                if (low[to] >= discovery[vertex])
                    extract_component(edge_id);
            } else if (discovery[to] < discovery[vertex]) {
                edge_stack.push_back({vertex, to, edge_id});
                low[vertex] = min(low[vertex], discovery[to]);
            }
        }
    }
};
\end{lstlisting}

Ejemplo de uso:

\begin{lstlisting}[style=cpp]
vector<vector<pair<int,int>>> graph(n);
for (int id = 0; id < m; ++id) {
    int u, v;
    cin >> u >> v;
    graph[u].push_back({v, id});
    graph[v].push_back({u, id});
}
VertexBiconnectedComponents bcc(graph);
for (const auto &component : bcc.components) {
    for (int vertex : component) cout << vertex << ' ';
    cout << '\n';
}
\end{lstlisting}

Un vértice es de articulación si aparece en más de un bloque:

\begin{lstlisting}[style=cpp]
vector<int> appearances(n);
for (const auto &component : bcc.components)
    for (int vertex : component)
        ++appearances[vertex];

for (int vertex = 0; vertex < n; ++vertex)
    if (appearances[vertex] > 1)
        cout << vertex << " es articulacion\n";
\end{lstlisting}

\warning{Cada arista no dirigida debe usar el mismo \texttt{edge\_id} en ambas
direcciones. Los vértices de articulación aparecen en varias componentes; esto
es intencional.}